% ============================================================================
% Project WILSON — Αναλυτική Τεχνική Αναφορά
% ----------------------------------------------------------------------------
% Πλήρης αναφορά με έμφαση στην υπολογιστική όραση (drone pipeline) και την
% τρισδιάστατη γεωμετρία (γεωαναφορά pixel->GPS, ανακατασκευή αναγλύφου,
% πεδίο ανέμου). Μεταγλώττιση με XeLaTeX.
%
% ΣΧΗΜΑΤΑ: όλα τα \benchfig δείχνουν σε report/figures/<name>.png.
% Όσο τα αρχεία λείπουν, τυπώνεται πλαίσιο-placeholder ώστε το έγγραφο να
% μεταγλωττίζεται κανονικά. Δες Παράρτημα Β για το πώς παράγεται κάθε σχήμα.
% ============================================================================
\documentclass[11pt]{article}
\usepackage[margin=2.3cm]{geometry}
\usepackage{fontspec}
% DejaVu όπου υπάρχει (Linux/CI), αλλιώς fallback σε fonts συστήματος macOS
% με πλήρη ελληνική κάλυψη.
\IfFontExistsTF{DejaVu Serif}
  {\setmainfont{DejaVu Serif}%
   \setsansfont{DejaVu Sans}%
   \setmonofont{DejaVu Sans Mono}[Scale=0.86]}
  {\setmainfont{Georgia}%
   \setsansfont{Arial}%
   \setmonofont{Menlo}[Scale=0.86]}
\usepackage{amsmath, amssymb}
\usepackage{graphicx}
\graphicspath{{figures/}}
\usepackage{booktabs}
\usepackage{enumitem}
\setlist{nosep, leftmargin=1.4em}
\usepackage{hyperref}
\hypersetup{hidelinks}
\setcounter{tocdepth}{2}

\newcommand{\code}[1]{\texttt{#1}}

% ── Benchmark figure με placeholder όταν λείπει η εικόνα ────────────────────
% \benchfig[πλάτος]{αρχείο.png}{λεζάντα}{label}
\newcommand{\benchfig}[4][0.82\linewidth]{%
\begin{figure}[htbp]\centering
\IfFileExists{figures/#2}{\includegraphics[width=#1]{#2}}{%
\fbox{\parbox[c][0.32\linewidth][c]{#1}{\centering
\ttfamily [ΕΚΚΡΕΜΕΙ ΕΙΚΟΝΑ]\\ \detokenize{figures/#2}\\[4pt]
\footnotesize βλ.\ Παράρτημα Β για την εντολή παραγωγής}}}%
\caption{#3}\label{#4}
\end{figure}}

\title{\textbf{Project WILSON}\\[2pt]
\large Διαδραστικός Προσομοιωτής Εξάπλωσης Δασικών Πυρκαγιών
με Βαθμονόμηση σε Πραγματικά Ελληνικά Συμβάντα\\[4pt]
\normalsize Αναλυτική τεχνική αναφορά --- με έμφαση στην υπολογιστική όραση
και την τρισδιάστατη γεωμετρία}
\date{Ιούλιος 2026}

\begin{document}
\maketitle
\tableofcontents
\newpage

% ============================================================================
\section{Εισαγωγή και στόχος του συστήματος}

Το WILSON είναι ένα πλήρες σύστημα προσομοίωσης εξάπλωσης δασικής πυρκαγιάς
σε πραγματικό ελληνικό ανάγλυφο. Συνδυάζει: (α) φυσική Rothermel πάνω σε
κυψελωτό αυτόματο (cellular automaton), (β) πραγματικά δεδομένα εδάφους,
καύσιμης ύλης και καιρού, (γ) διαδραστικό χάρτη στον browser με εργαλεία
επέμβασης (αντιπυρικές ζώνες, ρίψεις νερού, γραμμές ανάσχεσης), (δ) δύο
ανεξάρτητους βελτιστοποιητές βαθμονόμησης (Particle Swarm Optimization και
Differential Evolution), (ε) επικύρωση (hindcast) έναντι δορυφορικών
περιμέτρων πραγματικών πυρκαγιών (Εύβοια 2021, Ρόδος 2023, Έβρος 2023),
(στ) ένα πειραματικό υποσύστημα \emph{υπολογιστικής όρασης από drone} που
εντοπίζει ενεργή φλόγα σε βίντεο και τη γεωαναφέρει σε συντεταγμένες WGS-84
μέσω μοντέλου pinhole κάμερας, και (ζ) τρισδιάστατη απεικόνιση αναγλύφου με
διορθωμένο από το έδαφος πεδίο ανέμου.

Αρχιτεκτονικά, το σύστημα αποτελείται από Python backend (WebSocket server,
port 8765), web client (MapLibre, \code{sandbox.html}), βοηθητικό Flask API
για τρισδιάστατα πλέγματα (\code{mesh\_api.py}, port 5000) και το αυτόνομο
drone pipeline (\code{drone/src/}). Η παρούσα αναφορά περιγράφει όλα τα
υποσυστήματα, με ιδιαίτερο βάθος στα κεφάλαια~\ref{sec:cv}
(υπολογιστική όραση) και~\ref{sec:geom} (τρισδιάστατη γεωμετρία), και
παραθέτει μετρήσεις επιδόσεων· ο πλήρης πίνακας δεδομένων επικύρωσης
βρίσκεται στο συνοδευτικό \code{report\_data.tex}.

% ============================================================================
\section{Πυρήνας φυσικής προσομοίωσης (\code{core/})}

\subsection{Μοντέλο φωτιάς \code{core/fire\_model.py}}

Ο πυρήνας είναι η κλάση \code{CellularAutomataFire} ($\sim$900 γραμμές,
πλήρως vectorised NumPy). Κάθε κελί του πλέγματος βρίσκεται σε μία από
τρεις καταστάσεις (0 άκαυτο, 1 καιγόμενο, 2 καμένο) και συσσωρεύει
θερμότητα στο πεδίο \code{ignition\_fraction}· όταν αυτό ξεπεράσει το
ρυθμιζόμενο κατώφλι ανάφλεξης (\code{\_ignition\_threshold}), το κελί
αναφλέγεται και καίει για \code{\_burn\_time\_steps} βήματα (κλιμακούμενα
με το χρονικό βήμα $dt$).

Η μέθοδος \code{\_precompute\_ros\_grid()} υλοποιεί την πλήρη εξίσωση
Rothermel (1972) ανά κελί --- ένταση αντίδρασης (reaction intensity),
λόγος συμπίεσης (packing ratio), απόσβεση υγρασίας $\eta_m$, παράγοντας
ανέμου $\phi_w$ (με όριο 1500~ft/min) και παράγοντας κλίσης $\phi_s$ ---
και παράγει πίνακα πιθανοτήτων εξάπλωσης \code{p\_spread} διαστάσεων
$8\times R\times C$ (8 γείτονες). Σημαντική λεπτομέρεια υλοποίησης: κάθε
ανακατασκευή του πίνακα (π.χ.\ σε ωριαία ενημέρωση καιρού) μηδενίζει
εξωτερικές τροποποιήσεις, γεγονός που επηρέασε τη σχεδίαση της
βαθμονόμησης (\S\ref{sec:calib}).

Επιπλέον μηχανισμοί:
\begin{itemize}
  \item \emph{Monte-Carlo spotting} (τύπου Albini): πιθανότητα εκτόξευσης
        καύτρας ανάλογη της ταχύτητας ανέμου, log-normal κατανομή απόστασης
        πτήσης, απόσβεση από την τοπική ισχύ γραμμών ανάσχεσης στο σημείο
        πτώσης.
  \item \emph{Νερό/κατάσβεση}: πιθανοτική κατάσβεση με σιγμοειδή συνάρτηση
        (ποσότητα νερού έναντι φορτίου καύσιμης ύλης, υγρασίας, ανέμου),
        περιορισμένη στο bounding box των καιγόμενων κελιών για επιδόσεις.
  \item \emph{Ντετερμινισμός}: σταθερό RNG seed για αναπαραγωγιμότητα.
\end{itemize}

Επιδόσεις ανά βήμα προσομοίωσης: 0{,}9~ms σε πλέγμα $200^2$, 3{,}5~ms σε
$400^2$, 16~ms σε $800^2$.

\benchfig{ros_curves.png}{Χαρακτηρισμός φυσικής: ταχύτητα εξάπλωσης (ROS)
συναρτήσει ταχύτητας ανέμου, κλίσης και υγρασίας καύσιμης ύλης, ανά
αντιπροσωπευτική κλάση βλάστησης. Το αναλλοίωτο ανάλυσης (100$\to$50~m)
παραμένει εντός 3\%.}{fig:ros-curves}

\subsection{Τοπίο και καύσιμη ύλη \code{core/landscape.py}, \code{core/fuels.py}}

Ορίζονται 23 ελληνικές κλάσεις καύσιμης ύλης (\code{GREEK\_FUELS}): πεύκα
(χαλέπιος, τραχεία κ.λπ.), μακία, φρύγανα, ξηρά χόρτα, ελαιώνες, δρυοδάση
κ.ά., με πλήρεις παραμέτρους Rothermel (fuel load, surface-area-to-volume
ratio, moisture of extinction, bulk density), συν 4 μη καύσιμες κλάσεις
(νερό, αστικός ιστός, δρόμοι, άγονο). Η \code{load\_real\_terrain()}
συνθέτει: DEM (GeoTIFF) $+$ CORINE land cover (PNG) χαρτογραφημένο σε
κλάσεις καύσιμης ύλης $+$ OSM overlay (δρόμοι $\to$ \code{Urban\_Road},
οικιστικά πολύγωνα $\to$ \code{Urban\_Fabric}) μέσω rasterio. Το πλέγμα
είναι south-up (γραμμή 0 $=$ νότος). Η υγρασία της καύσιμης ύλης προκύπτει
από το Equilibrium Moisture Content (θερμοκρασία / σχετική υγρασία).

\subsection{Πεδίο ανέμου \code{air/air.py}}\label{sec:wind}

Ο άνεμος midflame προκύπτει από τρία διαδοχικά στρώματα διόρθωσης πάνω
στον βασικό μετεωρολογικό άνεμο:
\begin{enumerate}
  \item \emph{Wind Adjustment Factor (WAF)} ανά κλάση βλάστησης: πυκνό
        δάσος $\approx 0{,}10\times$, χόρτα $\approx 0{,}40\times$
        (μετρημένο στο validation).
  \item \emph{Ορογραφική ενίσχυση}: επιτάχυνση σε ανηφορική ροή με
        πρόσθετο μπόνους νότιου προσανατολισμού.
  \item \emph{Mass-consistent διόρθωση} τύπου WindNinja: διαγνωστική
        επίλυση Poisson
        \[
          \nabla^2\phi = \mathrm{div}(U,V),\qquad
          (U',V') = (U,V) - \nabla\phi,
        \]
        με μετασχηματισμό DCT-II (συνοριακές συνθήκες Neumann,
        πολυπλοκότητα $O(NM\log NM)$), ώστε το τελικό πεδίο να είναι
        divergence-free: επιτάχυνση σε κορυφογραμμές, συμπίεση ροής σε
        κοιλάδες.
\end{enumerate}

\benchfig{wind_field.png}{Διορθωμένο πεδίο ανέμου πάνω σε πραγματικό
ανάγλυφο (streamlines): αριστερά ομοιόμορφος μετεωρολογικός άνεμος, δεξιά
μετά τα τρία στρώματα διόρθωσης (WAF, ορογραφία, mass-consistent Poisson).
Διακρίνεται η επιτάχυνση στις κορυφογραμμές.}{fig:windfield}

\subsection{Μικροκλιματική μνήμη \code{core/microclimate.py}}

Ο \code{MicroclimateLearner} διατηρεί rolling average χάρτη συχνότητας
καύσης \code{ignition\_freq} από προηγούμενα τρεξίματα, με σειριοποίηση
JSON ($\sim$1{,}4~MB, 25~ms, lossless round-trip), bilinear
\code{resize\_to()} μεταξύ αναλύσεων και \code{apply\_to\_model()} που
ενισχύει το \code{p\_spread} σε ιστορικά ενεργούς διαδρόμους. Ενημερώνεται
και από τα top-K αποτελέσματα του PSO. Μετρήθηκε ότι επιταχύνει την αρχική
σύγκλιση του PSO, αλλά μπορεί να μεροληπτεί την τελική βαθμονόμηση.

% ============================================================================
\section{Δεδομένα και pipelines απόκτησης (\code{pipeline/})}

\subsection{Αυτόματη λήψη δεδομένων \code{pipeline/auto\_fetcher.py}}
Ενιαίο σημείο λήψης με cache ανά bounding box στον δίσκο:
\begin{itemize}
  \item DEM: Copernicus COP30 μέσω terrain-tiles API (zoom 14), cache
        \code{terrain\_\{lat\}\_\{lon\}\_\{buf\}.tif}.
  \item CORINE land cover μέσω WMS σε PNG 200--800~px.
  \item OSM (Overpass): δρόμοι (motorway\ldots service) και οικιστικά
        πολύγωνα με πλήρη γεωμετρία (\code{out geom})· τρία mirrors με
        fallbacks.
  \item \code{fetch\_osm\_places()}: ονομασμένοι οικισμοί
        (place$=$city/town/village/hamlet/suburb) για τον Threat Advisor.
\end{itemize}

\subsection{Ιστορικές περίμετροι \code{pipeline/fired\_loader.py}}
Φόρτωση FIRED daily perimeters (GeoPackage 22~MB, 24.867 features, Ελλάδα
2000--2024). Γνωστά γεγονότα: Εύβοια 2021 (id 2871, 12 ημέρες, 52{,}5~kha),
Ρόδος 2023 (id 130111, 10 ημέρες, 19{,}9~kha), Έβρος 2023 (id 105620,
18 ημέρες, 89{,}4~kha). Η \code{get\_day1\_ignition\_seeds()} σπάει
MultiPolygon σε ξεχωριστά σημεία έναυσης (π.χ.\ Εύβοια: 2 εστίες)· η
περίμετρος rasterize-άρεται σε μάσκα πλέγματος για scoring.

\subsection{Καιρικά δεδομένα (ERA5)}
Οι \code{fetch\_weather / fetch\_weather\_hourly} (στο
\code{hindcast\_optimizer.py}) αντλούν από το Open-Meteo ERA5 archive:
άνεμος 10~m, διεύθυνση, θερμοκρασία, σχετική υγρασία ανά ώρα, με μετατροπή
της σύμβασης διεύθυνσης «από» $\to$ «προς». Εύρημα validation: το ERA5
0{,}25$^\circ$ υποεκτιμά συστηματικά τον άνεμο των ημερών έκρηξης κατά
0{,}4--3~m/s --- εξ ου και ο πολλαπλασιαστής ανέμου στη βαθμονόμηση.

% ============================================================================
\section{Βαθμονόμηση και αντίστροφη μοντελοποίηση}\label{sec:calib}

\subsection{Sandbox βελτιστοποιητής PSO \code{pipeline/particle\_optimizer.py}}
Particle Swarm 12 παραμέτρων που ταιριάζει την προσομοίωση σε χειροκίνητα
σχεδιασμένη περίμετρο «αληθείας» στο UI. Χώρος αναζήτησης: πολλαπλασιαστής
και στροφή ανέμου, offset υγρασίας, canopy WAF, slope factor, 5
πολλαπλασιαστές ROS ανά ομάδα καύσιμης ύλης, spotting rate, κατώφλι
ανάφλεξης. Τρέχει σε γρήγορο πλέγμα $200^2$ με ThreadPoolExecutor·
συνάρτηση καταλληλότητας το IoU στη χρονική στιγμή της μάσκας, συν
ensemble heatmap πρόβλεψης $+1$~h από τα top-50 τρεξίματα. Επιστρέφει και
το «παγωμένο» state του καλύτερου σωματιδίου, ώστε το «συνέχισε από το
best-IoU fire» να επαναφέρει \emph{ακριβώς} εκείνη τη φωτιά
(nearest-neighbour αναγωγή $200\to800$) αντί να ξανατρέχει στοχαστικά.

Παγίδα υλοποίησης: το \code{\_precompute\_ros\_grid()} σβήνει κάθε
εξωτερική τροποποίηση· ο canopy παράγοντας επιβιώνει πλέον ως attribute
(\code{\_canopy\_wind\_mult}) και οι λοιποί πολλαπλασιαστές
επανεφαρμόζονται μετά από κάθε rebuild (ωριαίος καιρός, ρίψεις νερού).

\benchfig{pso_convergence.png}{Σύγκλιση PSO: βέλτιστο IoU ανά γενιά για
5 ανεξάρτητα τρεξίματα, με και χωρίς προθέρμανση από τη μικροκλιματική
μνήμη.}{fig:pso-conv}

\subsection{Hindcast βελτιστοποιητής \code{pipeline/hindcast\_optimizer.py}}
Πλήρες end-to-end pipeline επικύρωσης σε πραγματικές πυρκαγιές: αλήθεια
FIRED/FIRMS, καιρός ERA5, πραγματικό έδαφος, βελτιστοποίηση Differential
Evolution σε 5 κρυφές μεταβλητές: offset υγρασίας $[-0{,}15, 0]$,
πολλαπλασιαστής ανέμου $[0{,}5, 6]$, στροφή ανέμου $\pm30^\circ$, κατώφλι
ανάφλεξης $[0{,}15, 1]$, πολλαπλασιαστής διάρκειας καύσης $[1, 5]$.
Σύνθετο σφάλμα:
\[
  E = 0{,}3\,(1-\mathrm{IoU}) + 0{,}7\,H_{\mathrm{norm}},
\]
όπου $H_{\mathrm{norm}}$ το συμμετρικό Hausdorff κανονικοποιημένο στη
διαγώνιο του πλέγματος --- δίνει ομαλή κλίση ακόμα και όταν οι μάσκες δεν
τέμνονται. Προεπιλογή σε FIRED mode είναι το \emph{perimeter assimilation}
(τύπου FARSITE validation): έναυση από τη σωρευτική παρατηρημένη περίμετρο
της ημέρας \code{ignition\_day} και scoring της σωρευτικής έκτασης στο
τέλος του παραθύρου· εναλλακτικά point-seed από centroids για «καθαρή»
from-ignition δεξιότητα. Η βαθμονόμηση γίνεται σε αδρό πλέγμα $200^2$
($dt=1$~min) και το τελικό IoU υπολογίζεται σε $800^2$ με CFL-σωστή
κλιμάκωση $dt\propto\Delta x$.

\subsection{Αποτελέσματα επικύρωσης (σύνοψη)}
Από 13 hindcast τρεξίματα: assimilated ημέρας-4 μέσο IoU $\approx 48\%$
(Εύβοια 53{,}2\%, Ρόδος 47{,}6\%, Έβρος 43{,}0\%)· ημέρας-6/8 έως
66{,}7--90{,}2\%. Point-seed ημέρας-1 μόνο $\sim$2\% όταν η ημέρα 1 είναι
σιγοκαίουσα (Εύβοια 43~ha, Έβρος 22~ha), αλλά 42{,}6\% στη Ρόδο (ημέρα 1
ήδη 968~ha). Συστηματική υπερπρόβλεψη $1{,}1$--$2{,}2\times$ (απουσία
μοντέλου κατάσβεσης, ομαλοποιημένος ERA5 άνεμος). Ο πλήρης πίνακας στο
\code{report\_data.tex}.

\benchfig{hindcast_evia.png}{Διαγνωστικό τετράπτυχο hindcast (Εύβοια 2021,
ημέρα 4, perimeter assimilation): παρατηρημένη περίμετρος FIRED,
προσομοιωμένη έκταση, επικάλυψη (IoU 53{,}2\%), χάρτης σφάλματος.
Παράγεται από την \code{plot\_results()}.}{fig:hindcast-quad}

% ============================================================================
\section{Υπολογιστική όραση: ανίχνευση φωτιάς από drone
(\code{drone/})}\label{sec:cv}

Το υποσύστημα drone υλοποιεί έναν \emph{υβριδικό} ανιχνευτή ενεργής φλόγας
σε ροή βίντεο: ένας βαθύς ανιχνευτής αντικειμένων (YOLOv8) εντοπίζει
\emph{πού} υπάρχει πιθανή φωτιά (bounding box), και στη συνέχεια κλασική
χρωματομετρική κατάτμηση οριοθετεί \emph{ακριβώς} τα pixels της φλόγας
μέσα στο box. Ο συνδυασμός δίνει και τη σημασιολογική στιβαρότητα του
νευρωνικού δικτύου (απόρριψη ήλιου, φώτων, πορτοκαλί αντικειμένων) και την
pixel-ακρίβεια της κατάτμησης, η οποία είναι απαραίτητη για τη γεωαναφορά
του κεφαλαίου~\ref{sec:geom}: το κέντρο βάρους του πολυγώνου φλόγας ---
όχι του χονδρικού box --- είναι αυτό που προβάλλεται στο έδαφος.

\benchfig{cv_pipeline.png}{Αρχιτεκτονική του pipeline όρασης: frame $\to$
YOLOv8 (πρόταση περιοχών) $\to$ χρωματική κατάτμηση Bravonoid (RGB) ή
θερμική κατωφλίωση (IR) $\to$ μορφολογία $\to$ πολύγωνο φλόγας, κέντρο
βάρους, ρυθμός οπτικής μεγέθυνσης.}{fig:cv-pipeline}

\subsection{Στάδιο 1: Ανιχνευτής YOLOv8 (\code{drone/src/vision.py})}

Χρησιμοποιείται μοντέλο YOLOv8-nano (\code{fire\_yolov8n.pt}, ultralytics)
με κατώφλι εμπιστοσύνης \code{conf=0.3}. Ο ανιχνευτής επιστρέφει boxes με
κλάσεις που ελέγχονται \emph{δυναμικά} από το \code{model.names}: κάθε
κλάση που περιέχει τη λέξη «fire» δρομολογείται στη λεπτομερή κατάτμηση,
ενώ οι υπόλοιπες (π.χ.\ «smoke») απλώς σημειώνονται στο πλαίσιο
προεπισκόπησης. Στην IR παραλλαγή (\code{drone/fire\_tracker.py}) η
συμπερασματολογία τρέχει σε \code{imgsz=320} για επιδόσεις edge. Τα boxes
ψαλιδίζονται στα όρια του frame πριν από την εξαγωγή crop.

Η επιλογή του «πρώτα ανιχνευτής, μετά κατάτμηση» έχει δύο πρακτικά
πλεονεκτήματα: (α) η ακριβή χρωματική ανάλυση τρέχει μόνο στο μικρό crop
του box και όχι σε ολόκληρο το frame, και (β) οι χρωματικοί κανόνες, που
μόνοι τους θα έδιναν ψευδώς θετικά σε δύση ηλίου ή κόκκινες στέγες,
εφαρμόζονται μόνο εκεί όπου το δίκτυο έχει ήδη δώσει σημασιολογική
ένδειξη φωτιάς.

\benchfig{yolo_detections.png}{Παραδείγματα ανίχνευσης σε πραγματικό
υλικό πτήσης (φάκελος \code{drone/data/vision\_check/}): επιτυχείς
ανιχνεύσεις με box (μπλε), πολύγωνο κατάτμησης (πράσινο) και κέντρο βάρους
(κόκκινο), μαζί με δύο αρνητικά frames χωρίς ανίχνευση για έλεγχο ψευδώς
θετικών.}{fig:yolo-frames}

\subsection{Στάδιο 2α: Χρωματομετρική κατάτμηση RGB/YCbCr («Bravonoid»)}

Μέσα στο crop της φωτιάς εφαρμόζεται συνδυασμός εμπειρικών χρωματικών
κανόνων σε δύο χρωματικούς χώρους.

\paragraph{Κανόνες RGB.} Για κάθε pixel με κανάλια $(R,G,B)$:
\begin{align*}
  \text{R1:}\quad & R > G > B \\
  \text{R2:}\quad & R > 190,\; G > 90,\; B < 140 \\
  \text{R3:}\quad & 0{,}1 \le \tfrac{G}{R+1} \le 1{,}0,\;\;
                    0{,}1 \le \tfrac{B}{R+1} \le 0{,}85,\;\;
                    0{,}1 \le \tfrac{B}{G+1} \le 0{,}85
\end{align*}
Το R1 κωδικοποιεί τη φασματική διάταξη της φλόγας (κόκκινο $>$ πράσινο $>$
μπλε), το R2 απαιτεί απόλυτη φωτεινότητα στο κόκκινο κανάλι, και το R3
περιορίζει τους λόγους καναλιών ώστε να αποκλείονται κορεσμένα λευκά και
γκρίζα (ο παρονομαστής $+1$ αποτρέπει διαίρεση με το μηδέν).

\paragraph{Κανόνες YCbCr.} Στον χώρο φωτεινότητας--χρωμικότητας:
\[
  \bigl(Y \ge 170 \;\lor\; Y < 145\bigr) \;\land\;
  50 \le C_b \le 120 \;\land\; 120 < C_r < 220 .
\]
Η φλόγα έχει χαρακτηριστικά υψηλό $C_r$ (κόκκινη χρωμικότητα) και χαμηλό
$C_b$· ο διαζευκτικός όρος στο $Y$ πιάνει τόσο τον υπέρλαμπρο πυρήνα όσο
και τις σκοτεινότερες παρυφές της φλόγας.

\paragraph{Ενίσχυση χαμηλού φωτισμού και συναίνεση.} Το crop περνά και από
παραλλαγή ενισχυμένη για χαμηλό φωτισμό: ισοστάθμιση ιστογράμματος στο
κανάλι $V$ του HSV (\code{equalizeHist}) και median blur $5\times5$. Οι
κανόνες εφαρμόζονται και στις δύο εκδοχές (κανονική/ενισχυμένη), δίνοντας
τέσσερις μάσκες: $\mathrm{RGB}$, $\mathrm{RGB}_{low}$, $\mathrm{YCbCr}$,
$\mathrm{YCbCr}_{low}$. Συνδυάζονται ως
\[
  M_{\cup} = (\mathrm{RGB}\lor \mathrm{YCbCr}) \land
             (\mathrm{RGB}_{low}\lor \mathrm{YCbCr}_{low}),
  \qquad
  M_{RGB} = \mathrm{RGB}\lor \mathrm{RGB}_{low},
\]
και επιλέγεται η $M_{RGB}$ όταν οι δύο μάσκες συμφωνούν σε $\ge 75\%$ των
pixels του crop (ένδειξη «εύκολης» σκηνής), αλλιώς η αυστηρότερη
$M_{\cup}$. Ο μηχανισμός αυτός λειτουργεί ως αυτόματη εναλλαγή μεταξύ
ευαισθησίας και ακρίβειας ανάλογα με τη δυσκολία της σκηνής.

\benchfig{bravonoid_masks.png}{Στάδια της κατάτμησης Bravonoid σε ένα crop
φωτιάς: αρχικό crop, ενισχυμένη low-light εκδοχή, μάσκα RGB, μάσκα YCbCr,
τελική μάσκα συναίνεσης, και το τελικό πολύγωνο μετά τη
μορφολογία.}{fig:bravonoid-masks}

\subsection{Στάδιο 2β: Θερμική (IR) κατάτμηση}

Για θερμικές κάμερες σε λειτουργία «White Hot»
(\code{drone/fire\_tracker.py}, \code{CAMERA\_MODE='IR'}) η χρωματική
ανάλυση αντικαθίσταται από κατωφλίωση έντασης: μετατροπή σε grayscale (αν
η κάμερα στέλνει pseudo-color παλέτα π.χ.\ Ironbow), Gaussian blur
$5\times5$ για τον έντονο θόρυβο των αισθητήρων μικροβολομέτρων, και
δυαδική κατωφλίωση στο 230/255 --- η φωτιά είναι απλώς τα θερμότερα
pixels. Το κατώφλι είναι ρυθμιζόμενο ανά κάμερα, αφού η αντιστοίχιση
θερμοκρασίας--έντασης εξαρτάται από το AGC της κάθε συσκευής.

\benchfig{ir_segmentation.png}{Θερμική κατάτμηση: IR frame (White Hot),
μάσκα κατωφλίου 230, τελικό πολύγωνο. Η IR οδός είναι αναίσθητη σε
συνθήκες φωτισμού και καπνό, αλλά απαιτεί βαθμονόμηση κατωφλίου ανά
αισθητήρα.}{fig:ir-seg}

\subsection{Στάδιο 3: Μορφολογία, πολύγωνο και οπτική τηλεμετρία}

Η δυαδική μάσκα καθαρίζεται με μορφολογικό κλείσιμο (ελλειπτικό στοιχείο
$4\times4$) και διπλή διαστολή, ώστε να γεφυρωθούν οι ασυνέχειες που
προκαλεί το τρεμόπαιγμα της φλόγας. Ακολουθεί \code{findContours}
(εξωτερικά περιγράμματα) και επιλέγεται το μεγαλύτερο κατά εμβαδόν, με
ελάχιστο αποδεκτό εμβαδόν 50~px$^2$ για απόρριψη θορύβου. Το περίγραμμα
μετατοπίζεται από τις συντεταγμένες του crop στις συντεταγμένες του
πλήρους frame και γίνεται το \emph{πολύγωνο φλόγας}.

Από το πολύγωνο υπολογίζονται:
\begin{itemize}
  \item \emph{Κέντρο βάρους} μέσω ροπών εικόνας:
        $c_x = m_{10}/m_{00}$, $c_y = m_{01}/m_{00}$ --- η είσοδος της
        γεωαναφοράς (\S\ref{sec:geom}).
  \item \emph{Ρυθμός οπτικής μεγέθυνσης} (optical growth rate): κυλιόμενο
        ιστορικό 15 δειγμάτων $(t_i, A_i)$ εμβαδού πολυγώνου και εκτίμηση
        $\dot A = \Delta A / \Delta t$ σε px/s μεταξύ πρώτου και
        τελευταίου δείγματος. Θετικός ρυθμός σηματοδοτεί ενεργά
        αναπτυσσόμενη εστία --- πρώιμο, αμιγώς οπτικό προειδοποιητικό
        σήμα πριν από κάθε φυσική μοντελοποίηση.
\end{itemize}

\benchfig{growth_rate.png}{Χρονοσειρά εμβαδού πολυγώνου φλόγας και
εξομαλυμένος ρυθμός μεγέθυνσης (px/s) σε δοκιμαστικό βίντεο. Οι κόκκινες
περιοχές αντιστοιχούν σε φάσεις ενεργής ανάπτυξης της
εστίας.}{fig:growth-rate}

\subsection{Επιδόσεις όρασης}

\benchfig{cv_fps.png}{Ρυθμός επεξεργασίας (FPS) του pipeline όρασης ανά
διαμόρφωση: YOLOv8n σε imgsz 320 έναντι 640, RGB (πλήρης Bravonoid)
έναντι IR (κατωφλίωση), σε CPU laptop-class. Η IR οδός είναι ελαφρύτερη
καθώς παρακάμπτει τους διπλούς χρωματικούς χώρους.}{fig:cv-fps}

Ο ανιχνευτής τρέχει σε πραγματικό χρόνο σε CPU χάρη στο μέγεθος nano του
YOLOv8 και στο ότι η κατάτμηση περιορίζεται στα crops. Το FPS εμφανίζεται
ζωντανά στο overlay του preview και τυπώνεται ανά frame στην κονσόλα.

% ============================================================================
\section{Τρισδιάστατη γεωμετρία Ι: γεωαναφορά pixel $\to$ GPS
(\code{drone/src/geometry.py})}\label{sec:geom}

Η καρδιά του drone pipeline είναι η προβολή ενός pixel φωτιάς $(u,v)$ σε
απόλυτη γεωγραφική συντεταγμένη WGS-84, με χρήση του μοντέλου pinhole
κάμερας και της ζωντανής τηλεμετρίας του drone. Η υλοποίηση
(\code{FireGeolocator.pixel\_to\_gps()}) ακολουθεί έξι βήματα με σαφώς
ορισμένα συστήματα συντεταγμένων.

\subsection{Συστήματα συντεταγμένων και συμβάσεις}

\begin{itemize}
  \item \textbf{NED} (North--East--Down): αδρανειακό πλαίσιο κόσμου,
        πρότυπο MAVLink/autopilot· ο άξονας $Z$ δείχνει προς τα
        \emph{κάτω}.
  \item \textbf{Body} (Forward--Right--Down): σωματόδετο πλαίσιο του
        drone.
  \item \textbf{Cam} (πλαίσιο OpenCV): $X$ δεξιά στην εικόνα, $Y$ κάτω,
        $Z$ κατά μήκος του οπτικού άξονα προς τη σκηνή.
\end{itemize}
Η προεπιλεγμένη τοποθέτηση κάμερας είναι \emph{nadir} (κατακόρυφα προς τα
κάτω), με τη σταθερή αντιστοίχιση
\[
  X_{cam}\!\to\!Y_{body},\quad Y_{cam}\!\to\!X_{body},\quad
  Z_{cam}\!\to\!Z_{body}
  \;\;\Longrightarrow\;\;
  R_{cam\to body}^{nadir} =
  \begin{pmatrix} 0&1&0\\ 1&0&0\\ 0&0&1 \end{pmatrix}.
\]
Για gimbal με κλίση εμπρός, μια πρόσθετη στροφή
$R_y(-\theta_{off})$ προ-συντίθεται μία φορά κατά την αρχικοποίηση
(\code{cam\_pitch\_offset\_deg}).

\benchfig{coordinate_frames.png}{Τα τρία συστήματα συντεταγμένων (NED,
Body, Cam) και οι μεταξύ τους στροφές. Σχήμα-ορισμός για όλες τις
εξισώσεις της ενότητας.}{fig:frames}

\subsection{Βήμα 1: Αντιστροφή του μοντέλου pinhole}

Με εσωτερικές παραμέτρους (intrinsics) $f_x, f_y$ (εστιακή απόσταση σε
pixels) και $c_x, c_y$ (κύριο σημείο), το pixel $(u,v)$ απεικονίζεται σε
κανονικοποιημένη διεύθυνση στο πλαίσιο κάμερας:
\[
  x_n = \frac{u - c_x}{f_x},\qquad
  y_n = \frac{v - c_y}{f_y},\qquad
  \mathbf d_{cam} = (x_n,\, y_n,\, 1)^\top .
\]
Το διάνυσμα δεν χρειάζεται κανονικοποίηση μέτρου, αφού στη συνέχεια
κλιμακώνεται ελεύθερα από την τομή με το έδαφος.

Για την κάμερα του DJI Spark (live feed 1280$\times$720, διαγώνιο FOV
81{,}9$^\circ$, αισθητήρας 1/2.3$''$) η εστιακή απόσταση εκτιμάται
αναλυτικά:
\[
  f \;=\; \frac{d_{px}/2}{\tan(\mathrm{FOV}/2)}
    \;=\; \frac{\sqrt{1280^2+720^2}/2}{\tan(40{,}95^\circ)}
    \;\approx\; 846\ \text{px},
\]
τιμή που χρησιμοποιείται στο προφίλ \code{spark} του \code{main.py}
($f_x{=}f_y{=}846$, $c_x{=}640$, $c_y{=}360$). Το γενικό προφίλ υποθέτει
1920$\times$1080 με $f\approx980$. Για εργασία ακριβείας προβλέπεται
πλήρης βαθμονόμηση με \code{cv2.calibrateCamera()} (και συντελεστές
παραμόρφωσης φακού, που προς το παρόν αγνοούνται --- βλ.\ περιορισμούς,
\S\ref{sec:limits}).

\subsection{Βήματα 2--3: Αλυσίδα στροφών cam $\to$ body $\to$ NED}

Το διάνυσμα περνά πρώτα στο σωματόδετο πλαίσιο,
$\mathbf d_{body} = R_{cam\to body}\,\mathbf d_{cam}$, και έπειτα στο NED
μέσω γωνιών Euler ZYX (yaw $\psi$ $\to$ pitch $\theta$ $\to$ roll $\phi$):
\[
  R_{body\to NED} = R_z(\psi)\, R_y(-\theta)\, R_x(-\phi),
  \qquad
  \mathbf d_{NED} = R_{body\to NED}\,\mathbf d_{body}.
\]
Τα αρνητικά πρόσημα στα $\theta,\phi$ κωδικοποιούν τη σύμβαση προσήμων
MAVLink (pitch-up θετικό, roll-right θετικό) σε σχέση με τη φορά
περιστροφής των στοιχειωδών πινάκων. Οι $R_x, R_y, R_z$ είναι οι
στοιχειώδεις πίνακες στροφής $3\times3$ (\code{\_Rx, \_Ry, \_Rz}).

\subsection{Βήμα 4: Τομή ακτίνας με το έδαφος (flat-earth)}

Στο NED το $Z$ είναι θετικό προς τα κάτω, οπότε το επίπεδο του εδάφους
βρίσκεται σε $z = h$, όπου $h$ το ύψος πτήσης πάνω από το έδαφος (AGL).
Η ακτίνα από το drone $\mathbf P(t) = t\,\mathbf d_{NED}$ τέμνει το έδαφος
όταν $t\, d_z = h$:
\[
  t = \frac{h}{d_z},\qquad
  \Delta N = t\, d_N,\qquad
  \Delta E = t\, d_E,\qquad
  \rho = \sqrt{\Delta N^2 + \Delta E^2},
\]
με $\rho$ την οριζόντια απόσταση drone--φωτιά. Αν $d_z \le 10^{-6}$
(ακτίνα οριζόντια ή προς τα πάνω, π.χ.\ σε ακραίο pitch), δεν υπάρχει
έγκυρη τομή και η συνάρτηση επιστρέφει \code{None}. Η υπόθεση επίπεδου
εδάφους είναι συνειδητή απλοποίηση· η ακριβής λύση θα έτεμνε την ακτίνα
με το DEM (βλ.\ \S\ref{sec:limits}).

\benchfig{ray_ground.png}{Γεωμετρία τομής ακτίνας--εδάφους: το pixel
$(u,v)$ ορίζει ακτίνα στο NED που τέμνει το επίπεδο $z=h$· σημειώνονται η
οριζόντια εμβέλεια $\rho$ και η γωνία πλαγιότητας της
ακτίνας.}{fig:raycast}

\subsection{Βήμα 5: Μετρικές μετατοπίσεις $\to$ WGS-84}

Οι μετατοπίσεις σε μέτρα μετατρέπονται σε γεωδαιτικές με equirectangular
προσέγγιση μικρών γωνιών (μέση ακτίνα Γης $R_E = 6\,371$~km):
\[
  \Delta\varphi = \frac{\Delta N}{R_E}\cdot\frac{180}{\pi},
  \qquad
  \Delta\lambda = \frac{\Delta E}{R_E \cos\varphi_0}\cdot\frac{180}{\pi},
\]
όπου $\varphi_0$ το γεωγραφικό πλάτος του drone. Για εμβέλειες εκατοντάδων
μέτρων το σφάλμα της προσέγγισης είναι αμελητέο ($\ll$ 1~m) σε σχέση με τις
αβεβαιότητες τηλεμετρίας.

\subsection{Βήμα 6: Ευρετική εμπιστοσύνη και ανάλυση σφάλματος}

Κάθε αποτέλεσμα συνοδεύεται από δείκτη εμπιστοσύνης
\[
  c = \mathrm{clip}\!\Bigl(1
      - \frac{\max(|\theta|,|\phi|)}{45^\circ}
      - \frac{\rho/h}{10},\; 0{,}05,\; 1\Bigr),
\]
που τιμωρεί (α) μεγάλες γωνίες στάσης (όπου τα σφάλματα τηλεμετρίας
μεγεθύνονται) και (β) πολύ πλάγιες ακτίνες (μεγάλο $\rho/h$).

Η θεωρητική ευαισθησία δικαιολογεί τη μορφή του δείκτη: για γωνία
πλαγιότητας $\alpha$ από το ναδίρ, η οριζόντια εμβέλεια είναι
$\rho = h\tan\alpha$, οπότε ένα σφάλμα στάσης $\delta\alpha$ (rad)
προκαλεί σφάλμα θέσης
\[
  \delta\rho = \frac{h}{\cos^2\alpha}\,\delta\alpha ,
\]
δηλαδή στο ναδίρ ($\alpha{=}0$) 1$^\circ$ σφάλματος gimbal σε ύψος 100~m
δίνει $\approx$1{,}7~m, ενώ σε $\alpha{=}60^\circ$ το ίδιο σφάλμα δίνει
$\approx$7~m --- τετραπλασιασμός. Αντίστοιχα, το χωρικό μέγεθος ενός pixel
στο έδαφος (GSD) στο ναδίρ είναι $h/f$ ($\approx$0{,}12~m/px για το Spark
στα 100~m), και το βαρομετρικό σφάλμα ύψους περνά αναλογικά στην εμβέλεια:
$\delta\rho/\rho = \delta h / h$.

\benchfig{geoloc_sensitivity.png}{Ευαισθησία γεωαναφοράς: αναμενόμενο
σφάλμα θέσης στο έδαφος συναρτήσει γωνίας πλαγιότητας της ακτίνας, για
σφάλματα στάσης 0{,}5$^\circ$/1$^\circ$/2$^\circ$ και ύψη 50/100/150~m
(αναλυτικές καμπύλες $h\,\delta\alpha/\cos^2\alpha$). Επάνω δεξιά: ο
ευρετικός δείκτης εμπιστοσύνης της υλοποίησης στον ίδιο
χώρο.}{fig:geoloc-sens}

\subsection{Τηλεμετρία DJI: από flight record σε γωνίες κάμερας
(\code{drone/src/dji\_log.py})}

Επειδή το DJI Spark δεν εκθέτει MAVLink, το pipeline υποστηρίζει offline
λειτουργία: το cached βίντεο του DJI GO 4 συγχρονίζεται με το flight
record (\code{.txt} $\to$ CSV μέσω dji-log-parser). Ο \code{DJILogBook}
χειρίζεται τις ιδιαιτερότητες των exporters:
\begin{itemize}
  \item \emph{Ασαφής αντιστοίχιση στηλών}: τα ονόματα κεφαλίδων διαφέρουν
        ανά exporter, οπότε κανονικοποιούνται (πεζά, μόνο αλφαριθμητικά)
        και ταιριάζονται σε λίστες υποψηφίων ανά λογικό πεδίο.
  \item \emph{Gimbal αντί σκάφους}: το gimbal σταθεροποιεί την κάμερα
        ανεξάρτητα από το σώμα, οπότε στη γεωμετρία τροφοδοτείται η στάση
        του \emph{gimbal}. Η σύμβαση DJI (0$^\circ$ οριζόντια, $-90^\circ$
        κατακόρυφα κάτω) απεικονίζεται στη σύμβαση nadir της γεωμετρίας ως
        $\theta_{telem} = -(90^\circ + \theta_{gimbal})$, ώστε gimbal
        $-90^\circ \to$ ναδίρ και gimbal $0^\circ \to$ ορίζοντας. Το
        Spark δεν έχει άξονα yaw στο gimbal, άρα heading κάμερας $=$ yaw
        σκάφους.
  \item \emph{Ξετύλιγμα yaw}: το yaw ξετυλίγεται (unwrap) πριν τη γραμμική
        παρεμβολή, ώστε η παρεμβολή να μην σαρώνει λανθασμένα το άλμα
        $360^\circ\to 0^\circ$.
  \item \emph{Συγχρονισμός βίντεο--log}: το frame 0 αγκυρώνεται στη στιγμή
        που άναψε η σημαία \code{isVideo} (έναρξη εγγραφής), με
        χειροκίνητη διόρθωση \code{--video-offset} σε δευτερόλεπτα· κάθε
        frame $k$ αντιστοιχίζεται σε χρόνο log $t_0 + k/\mathrm{fps}$ και
        η τηλεμετρία παρεμβάλλεται γραμμικά.
  \item \emph{Έλεγχοι υγιεινής}: αυτόματη ανίχνευση GPS σε ακτίνια
        (η Ελλάδα θα διαβαζόταν ως $\varphi\approx0{,}66$), απόρριψη
        γραμμών πριν το GPS lock (0,0), ταξινόμηση κατά χρόνο. Το
        βαρομετρικό ύψος πάνω από το σημείο απογείωσης χρησιμοποιείται ως
        AGL --- έγκυρο μόνο σε σχετικά επίπεδο έδαφος γύρω από το σημείο
        εκτόξευσης.
\end{itemize}

\benchfig{dji_telemetry.png}{Τηλεμετρία flight record (πτήση 2026-07-08):
ύψος AGL, gimbal pitch και yaw συναρτήσει χρόνου, με σκιασμένα τα
διαστήματα ενεργής εγγραφής βίντεο (\code{isVideo}) στα οποία αγκυρώνεται
ο συγχρονισμός frame--τηλεμετρίας.}{fig:dji-sync}

\subsection{Ενορχήστρωση και καταγραφή (\code{drone/src/main.py})}

Ο «ελεγκτής κυκλοφορίας» ενώνει όραση και γεωμετρία πίσω από μια
αφηρημένη διεπαφή \code{DroneInterface} με τρεις υλοποιήσεις: (1)
\code{MAVLinkInterface} για ArduPilot/PX4 (τηλεμετρία από
\code{GLOBAL\_POSITION\_INT}+\code{ATTITUDE} στα 10~Hz, μη-μπλοκάρουσα
αποστράγγιση μηνυμάτων, βίντεο από USB/RTSP), (2) \code{DJILogInterface}
(offline replay όπως παραπάνω), (3) \code{VideoFileInterface} με mock
τηλεμετρία για ανάπτυξη. Κάθε ανίχνευση καταγράφεται στο
\code{data/fire\_gps\_log.csv} (χρονοσφραγίδα UTC, θέση/στάση drone,
εκτιμώμενη θέση φωτιάς, εμβέλεια, εμπιστοσύνη, pixel), ώστε οι εντοπισμοί
να μπορούν μελλοντικά να τροφοδοτούν αυτόματες εναύσεις στον προσομοιωτή.

\benchfig{fire_gps_map.png}{Γεωαναφερμένες ανιχνεύσεις από το
\code{fire\_gps\_log.csv} πάνω σε δορυφορικό υπόβαθρο: ίχνος πτήσης του
drone (μπλε), εκτιμημένες θέσεις φωτιάς χρωματισμένες κατά δείκτη
εμπιστοσύνης, και η διασπορά της συστάδας εκτιμήσεων γύρω από την
πραγματική εστία.}{fig:gps-map}

% ============================================================================
\section{Τρισδιάστατη γεωμετρία ΙΙ: ανακατασκευή και απεικόνιση αναγλύφου}
\label{sec:mesh}

\subsection{Flask API πλεγμάτων \code{mesh\_api.py}}

Το endpoint \code{POST /mesh} δέχεται $\{$lat, lon, bbox, wind$\}$ και
επιστρέφει τρία διαδραστικά Plotly figures (JSON) συν πλέγμα ανέμου:

\begin{itemize}
  \item \emph{Γεωμετρικά σωστό αποτύπωμα}: η περιοχή είναι \emph{φυσικά}
        τετράγωνη 2$\times$2~km --- το γεωγραφικό μήκος διορθώνεται κατά
        $1/\cos\varphi$ ώστε πλάτος και ύψος σε μέτρα να ταυτίζονται
        (\code{\_square\_bounds()}), αντί για «τετράγωνο» σε μοίρες που
        στην Ελλάδα θα ήταν $\sim$21\% πλατύτερο απ' ό,τι ψηλό.
  \item \emph{Πηγές δεδομένων μέσω slippy tiles}: μετατροπή
        (lat,lon)$\to$(x,y,z) με τη Web-Mercator σχέση
        $y = \bigl(1 - \operatorname{asinh}(\tan\varphi)/\pi\bigr)2^{z-1}$·
        DEM από Terrarium tiles σε zoom 14 ($\approx$9{,}5~m/px) και
        δορυφορική υφή ESRI σε zoom 17 ($\approx$1{,}2~m/px).
  \item \emph{Τρεις αναπαραστάσεις}: (α) DEM surface με ντυμένη δορυφορική
        υφή, (β) νέφος σημείων (έως 4.000 σημεία), (γ) TIN
        (τριγωνοποίηση Delaunay) --- και οι τρεις σε μετρικούς άξονες
        X/Y/Z με aspect ratio $z=0{,}35$ για ρεαλιστική υπερύψωση.
  \item \emph{Βέλη ανέμου}: πλέγμα 10$\times$10 διανυσμάτων από το
        \code{air/air.py} (πλήρης αλυσίδα διόρθωσης της \S\ref{sec:wind})
        όταν δίνεται ταχύτητα/διεύθυνση, αλλιώς ομοιόμορφο fallback.
  \item Συνθετικό fallback DEM όταν αποτύχει η λήψη πραγματικού.
\end{itemize}

\benchfig{mesh_dem.png}{Plotly DEM surface 2$\times$2~km με δορυφορική
υφή ESRI z17 και βέλη διορθωμένου ανέμου (10$\times$10). Screenshot από
το 3D popup του UI.}{fig:mesh-dem}

\benchfig{mesh_tin.png}{Το ίδιο ανάγλυφο ως TIN (τριγωνοποίηση Delaunay)
και ως νέφος σημείων· οι τρεις αναπαραστάσεις επιστρέφονται μαζί από το
\code{POST /mesh}.}{fig:mesh-tin}

\subsection{Live 3D προβολή προσομοίωσης \code{simulation\_3d\_view.py}, \code{viz/}}

Το \code{simulation\_3d\_view.py} εξάγει την τρέχουσα περιοχή σε NPZ και
εκκινεί subprocess viewer με live ενημέρωση κατάστασης μέσω αρχείου: το
μέτωπο της φωτιάς αποδίδεται ως overlay πάνω στο τρισδιάστατο ανάγλυφο
όσο τρέχει η προσομοίωση. Ο φάκελος \code{viz/} περιέχει τους
standalone visualizers (2D και 3D).

% ============================================================================
\section{Server προσομοίωσης \code{server.py}}

WebSocket server (port 8765) που ενοποιεί όλα τα υποσυστήματα:
\begin{itemize}
  \item \emph{Πρωτόκολλο}: \code{init} (περιοχή, εναύσεις, καιρός,
        προαιρετικό optimizer mode) $\to$ \code{init\_ack} $+$ ροή
        \code{frame} (RLE GeoJSON καιγόμενων/καμένων· 7{,}3~ms $+$ 104~KB
        ανά frame σε $800^2$)· ξεχωριστά sessions \code{init\_optimizer}
        και \code{init\_hindcast}.
  \item \emph{Επεμβάσεις}: containment line (βαθμωτή απόσβεση $+$ νερό $+$
        υγρασία, με φθορά στον χρόνο), firebreak (σκληρό φράγμα που
        μπλοκάρει και τις καύτρες στη διαδρομή), water drop/brush
        (Monte-Carlo σταγόνες).
  \item \emph{Ωριαίος καιρός}: ενημέρωση από ERA5 όταν δίνεται ημερομηνία,
        με επανεφαρμογή της βαθμονόμησης στο \code{p\_spread} μετά από
        κάθε rebuild.
  \item \emph{Optimizer continue-modes}: \code{ground\_truth} (η
        περίμετρος της μάσκας καίγεται, το εσωτερικό καμένο) και
        \code{best\_params\_sim} (επαναφορά του παγωμένου best fire)· και
        τα δύο ξεκινούν \emph{παγωμένα} στην ηλικία της μάσκας.
  \item Ιστορικό snapshots για scrubbing, cell-explanation
        (Alt$+$click, πλήρης ανάλυση Rothermel ανά κελί), rect-analysis,
        ενημέρωση microclimate στο τέλος.
\end{itemize}

% ============================================================================
\section{Web διεπαφή \code{sandbox.html}}

Ενιαία σελίδα ($\sim$5.500 γραμμές) με χάρτη MapLibre και αριστερό rail
με συρτάρια: Setup \& Weather, Tools (εναύσεις, μάσκα αληθείας, νερό,
containment, firebreak ζώνες με save/load --- \code{firebreak.py}),
Optimizer (PSO $+$ συνέχιση προσομοίωσης), Hindcast Validation (3
προκαθορισμένες πυρκαγιές), Microclimate Memory, Fire Events, Threat
Advisor. Η απόδοση της φωτιάς συνδυάζει GeoJSON layers με canvas
glow/καπνό (particles), wind streamlines από το πεδίο του server, ηχητικό
feedback, θέμα light/dark (Catppuccin) και 41 προγραμματιστικά παραγόμενα
PNG εικονίδια (\code{assets/gen\_icons.py}). Τα frames συγχωνεύονται σε
requestAnimationFrame· υπάρχει scrubber ιστορικού και 3D popup
(\S\ref{sec:mesh}).

% ============================================================================
\section{Σύμβουλος απειλών (Threat Advisor)
\code{pipeline/threat\_advisor.py}}

Live εκτίμηση κινδύνου ανά οικισμό:
\begin{itemize}
  \item Απόσταση από το ενεργό μέτωπο μέσω distance transform της μάσκας
        καύσης, EMA ταχύτητας προσέγγισης, ETA $=$ απόσταση/ταχύτητα·
        κλιμάκωση safe $\to$ watch ($\le$6~h) $\to$ warning ($\le$2~h ή
        1{,}5~km) $\to$ critical ($\le$30~min ή 500~m) $\to$ impact.
  \item Γλωσσικές αναφορές \emph{μόνο σε κλιμάκωση} $+$ προτάσεις: σημείο
        ρίψης νερού στο 35\% της διαδρομής μετώπου$\to$οικισμού, γραμμή
        ανάσχεσης κάθετη στον άξονα προσέγγισης.
  \item Διαδρομές διαφυγής: Dijkstra σε γράφο OSM ($\sim$100k κόμβοι στην
        Εύβοια) με ποινές εγγύτητας φωτιάς ($\le$300~m $=$ αποκλεισμένο)·
        έξοδοι $=$ κόμβοι στο εξωτερικό 4\% του χάρτη· ειλικρινές
        fallback «shelter in place» όταν όλοι οι δρόμοι κόβονται.
  \item Εκτέλεση κάθε $\ge$4~s wall σε worker thread πάνω σε snapshot του
        state (30--50~ms). Παγίδα: ελληνικά χωριά μοιράζονται ονόματα ---
        όλα τα κλειδιά είναι \code{name@lat,lon}.
\end{itemize}
Συμπληρωματικά, το \code{fire\_info\_panel.py} μετατρέπει την τηλεμετρία
της προσομοίωσης σε φυσική γλώσσα (ορόσημα έκτασης, κατεύθυνση/ταχύτητα
εξάπλωσης με δειγματοληψία εδάφους $\sim$150~m μπροστά από το μέτωπο,
αναφορές καιρού και επεμβάσεων).

% ============================================================================
\section{Εκκίνηση και βοηθητικά}

\code{main.py}: single-command εκκίνηση των τριών servers (HTTP 8000 /
WebSocket 8765 / Mesh 5000) και άνοιγμα browser. \code{gui\_launcher.py}:
Tkinter GUI του hindcast pipeline βήμα-βήμα. \code{firebreak.py}:
αποθήκευση/φόρτωση αντιπυρικών ζωνών σε CSV/JSON. \code{config.py}:
καθολικές σταθερές (πλέγμα, \code{BURN\_TIME\_STEPS}, καιρικές
προεπιλογές).

% ============================================================================
\section{Επικύρωση, περιορισμοί και μελλοντικές επεκτάσεις}\label{sec:limits}

\subsection{Τι έχει μετρηθεί}
Το πλήρες παράρτημα μετρήσεων βρίσκεται στο \code{report\_data.tex/.pdf}:
hindcast sweep 13 τρεξιμάτων, καμπύλες φυσικής, αποτελεσματικότητα
επεμβάσεων (αντιπυρική ζώνη 53--54\%, ρίψεις 46\%, γραμμές ανάσχεσης
6--11\% με πλήρη απόδοση της αποτυχίας στις καύτρες), σύγκλιση PSO,
στατιστικά έργου.

\subsection{Γνωστοί περιορισμοί}
\begin{itemize}
  \item \emph{Προσομοίωση}: απουσία μοντέλου ενεργητικής κατάσβεσης
        (συστηματική υπερπρόβλεψη), επιφανειακή μόνο φωτιά (το crown fire
        απορροφάται από πολλαπλασιαστές), ο ERA5 0{,}25$^\circ$
        ομαλοποιεί τον άνεμο, το cumulative scoring διογκώνει το IoU των
        όψιμων ημερών, υπερβολή spotting σε κελιά 25~m.
  \item \emph{Όραση}: τα βάρη YOLO δεν έχουν fine-tuned σε ελληνικό
        υλικό· οι χρωματικοί κανόνες έχουν σταθερά κατώφλια (ευαίσθητα σε
        white balance)· το IR κατώφλι 230 απαιτεί ρύθμιση ανά κάμερα· ο
        ρυθμός μεγέθυνσης είναι σε px/s, όχι m$^2$/s (θα απαιτούσε
        προβολή του πολυγώνου στο έδαφος).
  \item \emph{Γεωμετρία}: υπόθεση επίπεδου εδάφους (σε πλαγιές η τομή με
        DEM θα διόρθωνε συστηματικό σφάλμα εμβέλειας), αγνόηση
        παραμόρφωσης φακού, βαρομετρικό AGL έγκυρο μόνο κοντά στο σημείο
        απογείωσης, equirectangular προσέγγιση (αμελητέα στα μεγέθη μας).
\end{itemize}

\subsection{Πιθανές επεκτάσεις}
Ωριαίος καιρός εντός του hindcast παραθύρου· proxy κατάσβεσης· scoring
χωρίς το assimilated seed· πληθυσμός οικισμών στην ιεράρχηση του Threat
Advisor· \emph{σύνδεση των drone ανιχνεύσεων με αυτόματη έναυση στον
προσομοιωτή} (το \code{fire\_gps\_log.csv} έχει ήδη το απαιτούμενο
σχήμα)· τομή ακτίνας με DEM αντί για flat-earth· βαθμονόμηση κάμερας με
chessboard· μετατροπή του optical growth σε φυσικές μονάδες μέσω της
γεωαναφοράς όλου του πολυγώνου (όχι μόνο του κέντρου βάρους).

% ============================================================================
\appendix
\section{Παράρτημα Α: Πίνακες μετρήσεων}
Βλ.\ συνοδευτικό \code{report\_data.tex} (μεταγλωττίζεται αυτόνομα ή τα
tables του μπορούν να γίνουν \code{\textbackslash input} εδώ).

\section{Παράρτημα Β: Οδηγίες αναπαραγωγής σχημάτων}\label{app:figs}

Όλα τα σχήματα αναμένονται στον φάκελο \code{report/figures/} με τα
ονόματα του πίνακα. Μέχρι να τοποθετηθούν, το έγγραφο τυπώνει πλαίσια
placeholder.

\begin{table}[htbp]
\centering\small
\begin{tabular}{@{}lp{9.2cm}@{}}
\toprule
\textbf{Αρχείο} & \textbf{Τρόπος παραγωγής} \\
\midrule
\code{ros\_curves.png} & Σουίτα χαρακτηρισμού φυσικής (sweep ROS--άνεμος/κλίση/υγρασία)· τα δεδομένα υπάρχουν στο \code{report\_data.tex}, ενότητα C. \\
\code{wind\_field.png} & Screenshot των wind streamlines στο UI, ή απευθείας quiver plot από το \code{air/air.py} πριν/μετά τη διόρθωση. \\
\code{pso\_convergence.png} & Καταγραφή best-IoU ανά γενιά από το \code{particle\_optimizer.py} (5 seeds, με/χωρίς microclimate warm-start). \\
\code{hindcast\_evia.png} & \code{python gui\_launcher.py} ή απευθείας \code{hindcast\_optimizer.py} --- το τετράπτυχο PNG της \code{plot\_results()}. \\
\code{cv\_pipeline.png} & Σχεδιάγραμμα (draw.io/TikZ) βάσει \S\ref{sec:cv}. \\
\code{yolo\_detections.png} & Μωσαϊκό από υπάρχοντα frames στο \code{drone/data/vision\_check/} (frame\_*.png, nodetect\_*.png). \\
\code{bravonoid\_masks.png} & Μικρό script που καλεί \code{full\_bravonoid\_segmentation()} σε ένα crop και σώζει τα ενδιάμεσα στάδια δίπλα-δίπλα. \\
\code{ir\_segmentation.png} & \code{python drone/fire\_tracker.py} σε \code{test\_ir.mp4} με αποθήκευση σταδίων της \code{thermal\_segmentation()}. \\
\code{growth\_rate.png} & Plot της χρονοσειράς \code{area\_history} από τρέξιμο του \code{fire\_tracker.py}. \\
\code{cv\_fps.png} & Bar chart από τις τιμές FPS του overlay σε 4 διαμορφώσεις (320/640 $\times$ RGB/IR). \\
\code{coordinate\_frames.png}, \code{ray\_ground.png} & Σχεδιαγράμματα (TikZ/draw.io) βάσει \S\ref{sec:geom}. \\
\code{geoloc\_sensitivity.png} & Αναλυτικές καμπύλες $h\,\delta\alpha/\cos^2\alpha$ --- 20 γραμμές matplotlib, δεν απαιτεί πτήση. \\
\code{fire\_gps\_map.png} & Plot του \code{drone/data/fire\_gps\_log.csv} πάνω σε δορυφορικό υπόβαθρο (folium/contextily). \\
\code{dji\_telemetry.png} & Plot στηλών height/gimbal\,pitch/yaw/isVideo από το \code{ExportCSV\_2026-07-08\_[15-30-54].csv} μέσω \code{DJILogBook}. \\
\code{mesh\_dem.png}, \code{mesh\_tin.png} & Screenshots από το 3D popup του UI (\code{POST /mesh} στο port 5000). \\
\bottomrule
\end{tabular}
\caption{Αντιστοίχιση σχημάτων--εντολών παραγωγής.}
\end{table}

\end{document}
